\documentclass[12pt,a4paper]{article}
\usepackage{amsmath,amssymb,amsthm}
\usepackage{hyperref}
\usepackage{geometry}
\usepackage{enumitem}
\geometry{margin=2.5cm}

\newtheorem{theorem}{Theorem}[section]
\newtheorem{corollary}[theorem]{Corollary}
\newtheorem{proposition}[theorem]{Proposition}
\theoremstyle{definition}
\newtheorem{definition}[theorem]{Definition}
\theoremstyle{remark}
\newtheorem{remark}[theorem]{Remark}

\title{Muon and Tau Masses from Recognition Science:\\
$\varphi^{11}$ and $\varphi^{6}$ Generation Spacings}

\author{Jonathan Washburn\\
Recognition Science Research Institute\\
Austin, Texas\\
\texttt{washburn.jonathan@gmail.com}}

\date{March 2026}

\begin{document}
\maketitle

\begin{abstract}
We derive the muon and tau masses from the Recognition Science
$\varphi$-ladder.  The lepton mass hierarchy is purely geometric:
$m_\mu/m_e = \varphi^{11}$ and $m_\tau/m_\mu = \varphi^{6}$, where
the generation spacings $\tau(1) = 11$ and $\tau(2) = 17$ arise
from the $D = 3$ cube geometry (11 passive edges and 17 wallpaper
groups).  Lepton rungs: $r_e = 2$, $r_\mu = 13$, $r_\tau = 19$.
The mass formulas are $m_\mu = E_{\mathrm{coh}} \varphi^{13}$ and
$m_\tau = E_{\mathrm{coh}} \varphi^{19}$ in RS-native units, with
$E_{\mathrm{coh}} = \varphi^{-5}\,\mathrm{eV} \approx 0.090$~eV.
The predicted ratios ($m_\mu/m_e \approx 199$, $m_\tau/m_\mu
\approx 17.9$) agree with experiment ($206.8$, $16.8$) at the
$\sim 4\%$ level.
\end{abstract}

%% ============================================================
\section{Recognition Science Framework}
\label{sec:framework}
%% ============================================================

The lepton mass hierarchy derivation rests on the unique cost
functional and the $D=3$ cube geometry it selects.

\begin{definition}[RS Cost Functional]
For $x > 0$: $J(x) = \tfrac{1}{2}(x + x^{-1}) - 1$.
\end{definition}

\noindent\textbf{Axiom A1 (Normalization):} $J(1) = 0$.\\
\textbf{Axiom A2 (Recognition Composition Law, RCL):}
$J(xy)+J(x/y)=2J(x)J(y)+2J(x)+2J(y)$ for all $x,y>0$.\\
\textbf{Axiom A3 (Calibration):} $J''_{\log}(0) = 1$, where
$J_{\log}(t) := J(e^t)$.

\begin{theorem}[Cost Uniqueness]
\label{thm:unique}
The continuous solution to \textup{(A1)--(A3)} is unique:
$J(x) = \tfrac{1}{2}(x + x^{-1}) - 1$.
\end{theorem}

\begin{proof}[Proof sketch]
Set $H(t) = J_{\log}(t)+1$.  Axiom A2 transforms into the d'Alembert
equation $H(s+t)+H(s-t)=2H(s)H(t)$.  By the Acz\'el classification
theorem~\cite{Aczel1966}, the unique continuous solution with $H(0)=1$
and $H''(0)=1$ is $H(t)=\cosh t$, giving
$J(x)=\tfrac{1}{2}(x+x^{-1})-1$.
\end{proof}

\noindent The golden ratio $\varphi=(1+\sqrt{5})/2$ is forced by the
RCL self-similarity; $D=3$ is forced by $2^D=8$.  The lepton mass
hierarchy follows from the cube's combinatorial invariants:
$E_{\rm passive}=11$ (passive edges) determines the $e\to\mu$ spacing
and $W=17$ (wallpaper groups) determines the $\mu\to\tau$ spacing.
At integer approximation: $m_\mu/m_e\approx\varphi^{11}$ and
$m_\tau/m_\mu\approx\varphi^6$.  A companion paper (RS\_Electron\_Mass)
derives sub-ppm precision via non-integer steps
$S_{e\to\mu}\approx11.08$ and $S_{\mu\to\tau}\approx5.87$.

%% ============================================================
\section{Introduction}
\label{sec:intro}
%% ============================================================

The charged lepton mass hierarchy---electron, muon, tau---spans
five orders of magnitude.  In the Standard Model, these masses are
free Yukawa parameters with no theoretical explanation.  Recognition
Science (RS) derives the hierarchy from the $\varphi$-ladder, a
discrete cost structure tied to the $D = 3$ cube geometry.

The RS framework posits $J(x) = \tfrac{1}{2}(x + x^{-1}) - 1$ as
the cost functional, $\varphi = (1+\sqrt{5})/2 \approx 1.618$ as the
golden ratio, and $E_{\mathrm{coh}} = \varphi^{-5}\,\mathrm{eV}
\approx 0.090$~eV as the coherence energy.  Masses are placed on
rungs $r$ of the ladder: $m = E_{\mathrm{coh}} \varphi^r$.  The
generation spacings $\tau(1)$ and $\tau(2)$ between rungs arise
from cube invariants.

%% ============================================================
\section{Generation Spacings from Cube Geometry}
\label{sec:generation-spacings}
%% ============================================================

\begin{definition}[Generation Spacings]
\label{def:tau}
For the $D = 3$ cube:
\begin{enumerate}[label=(\roman*)]
  \item $\tau(1) = 11$: passive edges.  The cube has 12 edges; one
  is ``active'' (the tick direction), leaving $12 - 1 = 11$ passive
  edges.
  \item $\tau(2) = 17$: wallpaper groups of the cube.  The 17
  two-dimensional wallpaper groups classify the cube's symmetry
  projections.
\end{enumerate}
\end{definition}

\begin{remark}
The 8-tick epoch structure ($2^D = 8$ for $D = 3$) and the color
number $N_c = 3$ from face pairs are also fixed by the cube.  The
generation spacings $\tau(1)$, $\tau(2)$ are thus geometric, not
free parameters.
\end{remark}

%% ============================================================
\section{Muon and Tau Masses}
\label{sec:masses}
%% ============================================================

The electron sits at rung $r_e = 2$ (the minimal nontrivial rung).
The muon and tau rungs are obtained by adding the generation spacings.

\begin{proposition}[Muon Rung and Mass]
\label{thm:muon}
$r_\mu = r_e + \tau(1) = 2 + 11 = 13$.  Hence
\begin{equation}
  m_\mu = E_{\mathrm{coh}} \,\varphi^{13}.
\end{equation}
\end{proposition}

\begin{remark}
The assignment $r_\mu = 13$ is a structural consequence of the $D=3$
cube's passive-edge count.  The precise correction to the spacing
(making $\tau(1)$ slightly non-integer) is derived in the companion
paper RS\_Electron\_Mass via the fine-structure constant $\alpha$.
\end{remark}

\begin{proposition}[Tau Rung and Mass]
\label{thm:tau}
$r_\tau = r_e + \tau(2) = 2 + 17 = 19$ (wallpaper spacing).  Equivalently,
$r_\tau = r_\mu + (\tau(2) - \tau(1)) = 13 + 6 = 19$.  Hence
\begin{equation}
  m_\tau = E_{\mathrm{coh}} \,\varphi^{19}.
\end{equation}
\end{proposition}

\begin{corollary}[Mass Ratios]
\label{cor:ratios}
\begin{align}
  \frac{m_\mu}{m_e} &= \varphi^{11} \approx 199.0, \\
  \frac{m_\tau}{m_\mu} &= \varphi^{6} \approx 17.94.
\end{align}
\end{corollary}

%% ============================================================
\section{Comparison with Experiment}
\label{sec:comparison}
%% ============================================================

\begin{table}[h]
\centering
\caption{Lepton mass ratios: RS prediction vs.\ experiment~\cite{PDG}.}
\begin{tabular}{lcc}
\hline
Quantity & RS & Experiment \\
\hline
$m_\mu/m_e$ & $\varphi^{11} \approx 199$ & $206.768\,2830(46)$ \\
$m_\tau/m_\mu$ & $\varphi^{6} \approx 17.9$ & $16.817\pm 0.004$ \\
\hline
\end{tabular}
\end{table}

The $m_\mu/m_e$ prediction is $\sim 4\%$ low; the $m_\tau/m_\mu$
prediction is $\sim 6\%$ high.  These discrepancies are comparable
to the $\varphi$-ladder precision seen in other RS mass predictions.
Radiative corrections and scale-dependent definitions may account
for part of the difference.

%% ============================================================
\section{Predictions and Falsifiers}
\label{sec:predictions}
%% ============================================================

\begin{enumerate}[label=(\roman*)]
\item \textbf{Mass ratios}: $m_\mu/m_e \in [190, 210]$ and
$m_\tau/m_\mu \in [16.5, 19]$.  Falsifier: ratios outside these
bands at $> 5\sigma$ would challenge the $\varphi$-ladder assignment.

\item \textbf{Rung structure}: $r_e = 2$, $r_\mu = 13$, $r_\tau =
19$.  Falsifier: any lepton mass requiring a different rung
assignment inconsistent with $\tau(1) = 11$, $\tau(2) = 17$ would
falsify the geometric derivation.

\item \textbf{Fourth generation}: RS predicts exactly three charged
lepton generations.  Falsifier: discovery of a fourth light charged
lepton.
\end{enumerate}

%% ============================================================
\section{Conclusion}
\label{sec:conclusion}
%% ============================================================

Lepton masses are geometric in RS: generation spacings $\tau(1) =
11$ and $\tau(2) = 17$ from $D = 3$ cube geometry determine the
rungs $r_\mu = 13$, $r_\tau = 19$.  Mass ratios $m_\mu/m_e =
\varphi^{11}$ and $m_\tau/m_\mu = \varphi^{6}$ agree with
experiment at the $\sim 4\%$ level.

\begin{thebibliography}{9}
\bibitem{RS_Axioms}
J.~Washburn, ``The Algebra of Reality: A Recognition Science Derivation
of Physical Law,'' \emph{Axioms} \textbf{15}(2), 90 (2026).

\bibitem{Aczel1966}
J.~Acz\'el, \textit{Lectures on Functional Equations and Their Applications},
Academic Press, New York (1966).

\bibitem{PDG}
S.~Navas et al.\ (Particle Data Group), Phys.\ Rev.\ D \textbf{110},
030001 (2024).
\end{thebibliography}

\end{document}
